\section{N0--N3: what was built and what it decided}
\label{sec:evolution}

Every NOW item is pre-registered to end in a continue, change, or stop reading of an experiment
result, never in \enquote{the build finished}. Rule R13 of the plan says so: a NOW item that ends in
infrastructure without a gate reading is a failure of the plan (\repo{REORIENTATION.md}, \S22).
\evobs{} N0 and N1 read their gates directly from a code check. N2's registered gate reads a specific
number against a specific floor. N3's registered gate, the ΔAUROC test against a gold standard, has
not run, because no gold has been acquired. What N3 reports below are its own internal
proof-of-concept checks, not its program gate. \Cref{tab:now-verdicts}
summarizes all four.

\begin{table}[tbp]
  \centering
  \small
  \caption{N0--N3: purpose, what was built, and the registered verdict each one read. Test counts
  and check results are the live-run numbers used throughout this report
  (\cref{sec:architecture,sec:n2,sec:n3}). N3's key commit, \code{ec45c31}, is the N3 PoC commit;
  \code{754e2e6}, cited elsewhere in this report, is the later ADR-0008 commit.}
  \label{tab:now-verdicts}
  \begin{tabularx}{\linewidth}{@{}lXXXl@{}}
    \toprule
    Item & Purpose & What was built & Registered verdict & Key commit \\
    \midrule
    N0 & Reset after the reorientation: fix the DOI-checking hook, freeze Track A, create the
      domain-neutral package & \repo{scripts/hook_tests/test_check_dois.py}; \repo{residual/}
      scaffold with its own test hook & \textbf{Done, gate passed} (2026-09-28): the hook
      demonstrably fires on a test edit & \code{d4fbc60} \\
    N1 & Pre-register and build the evidence model: claim record, six epistemic labels, gap-map
      field types, a freeze & \repo{residual/src/residual/} (claims, vocab, ledger, provenance,
      gates, gapmap, freeze, residual); \repo{residual/frozen/n1.json} & \textbf{Done, gate:
      continue with 2 additions} (\code{KnowledgeType.INTERPRETATION},
      \code{ResponseDistribution}), at the registered ≤2 limit & \code{accb506} \\
    N2 & Build a gated reconstruction pipeline and test it live and against planted falsehoods
      and private questions & One-way pipeline (plan→search→fetch→extract→locate→verify→
      contradict→assemble); E-LIVE runs on PLC and GDPR; E-PLANT/E-ABST harnesses & \textbf{Engineering
      PoC: complete. Full research validation: deferred.} Registered gate: \textbf{INCONCLUSIVE}
      -- E-PLANT's Continue branch closed on the validity floor ($a_U{=}5/24 < 8$), and the gated
      arms of E-PLANT and E-ABST were never scored & \code{ea5c18e} \\
    N3 & Build a methodology-guided gap map over N2's ledgers that turns silence into candidate
      hypotheses and expert questions & Seven lenses over three ledgers (PLC, GDPR v1, GDPR v2);
      ranked maps; built-in adversarial checks (anti-renaming, closure null, mismatched-evidence
      control, cross-run stability) & \textbf{PoC: pipeline complete, hidden-knowledge prediction
      not demonstrated.} Closure step is \code{closure-uninformative} on every lens tested. Program
      gate (ΔAUROC vs.\ gold) \textbf{not run} -- gold not acquired & \code{ec45c31} \\
    \bottomrule
  \end{tabularx}
\end{table}

N0 and N1 close cleanly: each reads a code-checked condition (a hook firing; a field-type count
against a pre-registered ceiling) and both pass. \evobs{} N2's registered research verdict is
INCONCLUSIVE. The ungated E-PLANT baseline adopted only 5 of 24 planted falsehoods
($a_U = 5/24$, below the required 8). By the pre-registered validity-floor rule, this closes the
\emph{Continue} branch before the gated arms of E-PLANT and E-ABST were scored, and those arms are
the actual test of whether gating helps. The owner then chose not to spend further budget. Those arms
and a post-fix live re-check (E-LIVE v3) stay pre-registered and unexecuted. Separately, N2's
engineering milestone is complete (\cref{sec:architecture}).

N3 built and ran its full pipeline on real ledgers. Its mismatched-evidence-control check is built into
\code{gapmap/checks.py}, added after an adversarial review found the weakness, and re-run on every
replay. It shows that the closure judge does not beat the mismatched-evidence control on any lens
tested (lenses with fewer than two candidates are skipped). The own rate equals the control rate
($0.889/0.889$, $0.961/0.961$, $0.958/0.958$). The gap map is therefore a
working mechanism that produces \emph{candidates}, with no demonstrated ability yet to tell a real
absence from a topically nearest neighbor's evidence.

\subsection{The owner exception: N3 on N2's engineering milestone}\label{sec:owner-exception}

N2's own registered gate stayed INCONCLUSIVE. Under the plan's rule R13, that would block the next
item from starting. \evobs{} The owner made a single, recorded exception (\repo{REORIENTATION.md},
\S22; \repo{PROGRESS.md}). N2's engineering milestone was treated as sufficient to start N3; N2's
research validation stays deferred. N3's own stop rule is unaffected, since it reads N3's own
ΔAUROC test, not N2's gate.

\evint{} This was a softened gate, and the person who softened it is the person whose program depends
on it. No one outside the project read the gate or approved the exception. Recording the exception is
better than leaving it implicit, but it is not independent control. The program therefore adds an
external gate reader: an independent statistician, or a small advisory board with one, reads the
closure, screening and H2 gates (MS1, MS3, MS5) against the pre-registration and has authority over
them. Any later exception by the owner needs the reader's sign-off (\cref{sec:validation}).
